\documentclass[11pt,a4paper]{article}
\usepackage[utf8]{inputenc}
\usepackage[T1]{fontenc}
\usepackage[spanish,es-tabla]{babel}
\usepackage[margin=2.3cm]{geometry}
\usepackage{amsmath,amssymb}
\usepackage{booktabs}
\usepackage{longtable}
\usepackage{tabularx}
\usepackage{array}
\usepackage{xcolor}
\usepackage{colortbl}
\usepackage{hyperref}
\usepackage{fancyhdr}
\usepackage{titlesec}
\usepackage{graphicx}
\usepackage{listings}
\usepackage{tcolorbox}
\usepackage{enumitem}
\usepackage{multirow}
\tcbuselibrary{skins,breakable}

% Paleta de colores ejecutiva
\definecolor{primaryblue}{RGB}{16, 44, 87}
\definecolor{secondaryblue}{RGB}{53, 89, 143}
\definecolor{accentteal}{RGB}{20, 110, 120}
\definecolor{lightbg}{RGB}{245, 247, 250}
\definecolor{darkgray}{RGB}{50, 50, 50}
\definecolor{warnred}{RGB}{180, 30, 30}
\definecolor{okgreen}{RGB}{20, 110, 50}
\definecolor{warnorange}{RGB}{200, 100, 20}

\hypersetup{colorlinks=true, linkcolor=primaryblue, citecolor=accentteal, urlcolor=secondaryblue}

\pagestyle{fancy}
\fancyhf{}
\rhead{\textcolor{darkgray}{\small Entrega Final BI $|$ Persona 3: Arquitecto Dimensional RM}}
\lhead{\textcolor{darkgray}{\small Esquema Estrella, Grano, SCD Tipo 2 y Aditividad de Tasas/KPIs}}
\cfoot{\thepage}

\titleformat{\section}{\Large\bfseries\color{primaryblue}}{\thesection.}{0.5em}{}[\titlerule]
\titleformat{\subsection}{\large\bfseries\color{secondaryblue}}{\thesubsection.}{0.4em}{}
\titleformat{\subsubsection}{\normalsize\bfseries\color{accentteal}}{\thesubsubsection.}{0.3em}{}

\lstset{
  basicstyle=\ttfamily\footnotesize,
  backgroundcolor=\color{lightbg},
  frame=single,
  rulecolor=\color{gray!30},
  breaklines=true,
  tabsize=2,
  showstringspaces=false
}

\begin{document}

% Portada / Cabecera Ejecutiva
\begin{center}
    {\Huge\bfseries\color{primaryblue} Proyecto Final de Business Intelligence}\\[0.25cm]
    {\LARGE\bfseries\color{secondaryblue} DataMart y Esquema Estrella Enriquecido (2012--2026)}\\[0.2cm]
    {\large\bfseries\color{accentteal} Enfoque Estratégico 100\% Región Metropolitana de Santiago}\\[0.4cm]
    \rule{\linewidth}{1.2pt}\\[0.3cm]
    \textbf{Rol Evaluado:} Persona 3 --- Arquitecto Dimensional y Diccionario de Tasas/KPIs\\
    \textbf{Autor:} Angel Pino \quad $|$ \quad \textbf{Carrera:} Ingeniería en Negocios\\
    \textbf{Fecha de Entrega y Defensa Oficial:} 18 de Noviembre de 2026\\
    \textbf{Universo Auditado:} 58.926 Ofertas Académicas RM $|$ 69 Instituciones Activas en 2026 (SUBDERE 13)
\end{center}

\vspace{0.4cm}

\begin{tcolorbox}[colback=primaryblue!5!white,colframe=primaryblue,title=\textbf{Síntesis de Cumplimiento Técnico y Feedback Docente Aplicado}]
\begin{itemize}[leftmargin=*]
    \item \textbf{Foco Exclusivo en la Región Metropolitana (RM):} La tabla de hechos se acota estrictamente a los \textbf{58.926 registros históricos de Santiago} y las \textbf{69 instituciones activas en 2026} (35 Universidades, 19 IP, 11 CFT y 4 FF.AA.).
    \item \textbf{Tasas y KPIs en Tablas dim\_* y fact\_*:} Se superó la limitación de los conteos brutos. Cada tasa representa una razón explicativa que fundamenta decisiones de sobredemanda, capacidad de laboratorios (caso 120$\rightarrow$140), retención y selectividad PAES.
    \item \textbf{Cumplimiento Clase 4 DW (Pág. 62):} Cero columnas de texto \texttt{VARCHAR} en la tabla de hechos (100\% numérica) y eliminación absoluta de abreviaturas crudas en dimensiones (\texttt{'Universidad'}, \texttt{'Instituto Profesional'}, \texttt{'Centro de Formación Técnica'}, \texttt{'Institución de las Fuerzas Armadas y de Orden'}) y prefijos regulatorios \texttt{(a)/(b)}.
    \item \textbf{Integración de 2 Fuentes Externas:} Proyecciones demográficas del \textbf{INE RM} en \texttt{dim\_tiempo\_ine\_rm} y acreditación institucional \textbf{SIES/CNA} en \texttt{dim\_institucion\_sies}.
    \item \textbf{Certificación 1:N:} \textbf{0 registros huérfanos} verificados en las cuatro relaciones dimensionales sobre las 58.926 filas.
\end{itemize}
\end{tcolorbox}

\section{Topología Relacional: Esquema Estrella Enriquecido}

El diseño dimensional adoptado corresponde a un \textbf{Esquema Estrella Puro (Star Schema)}, compuesto por una tabla de hechos central (\texttt{fact\_matricula\_admision}) y cuatro tablas de dimensiones canónicas desnormalizadas, complementadas por una tabla agregada de optimización OLAP y una tabla puente de gobernanza:

\begin{table}[h!]
\centering
\small
\begin{tabularx}{\linewidth}{l c X}
\toprule
\textbf{Tabla / Hoja Excel} & \textbf{Filas (RM)} & \textbf{Rol y Justificación Arquitectónica (Clase 4 DW)} \\
\midrule
\texttt{fact\_matricula\_admision} & 58.926 & Tabla de hechos central (100\% numérica, cero columnas \texttt{VARCHAR}). \\
\texttt{dim\_tiempo\_ine\_rm} & 15 & Dimensión temporal (2012--2026) que integra proyecciones del \textbf{INE RM}. \\
\texttt{dim\_institucion\_sies} & 1.298 & Dimensión institucional con \textbf{SCD Tipo 2} y acreditación oficial \textbf{SIES}. \\
\texttt{dim\_sede_geografia} & 22 & Dimensión geográfica desnormalizada de sedes en la RM con tasas de ocupación. \\
\texttt{dim\_carrera\_programa} & 5.184 & Dimensión curricular con métricas de saturación sobre laboratorios base 35. \\
\texttt{agg\_resumen\_anual\_institucion\_rm} & 1.298 & Tabla agregada anual para consultas analíticas de alta velocidad (Pág. 63). \\
\texttt{puente\_institucion_sies} & 259 & Tabla puente oficial de homologación y auditoría CNED a SIES (Pauta Pág. 3). \\
\bottomrule
\end{tabularx}
\caption{Inventario Canónico de Tablas del DataMart de la Región Metropolitana.}
\end{table}

\subsection{Justificación de la Topología Estrella frente a Snowflake}
De acuerdo con los principios de Kimball y la cátedra de Data Warehousing (Clase 4, págs. 21--25), el Esquema Estrella minimiza el costo computacional de combinaciones relacionales (\emph{joins}), logrando que cualquier consulta multidimensional requiera a lo más un solo salto relacional por dimensión. Al desnormalizar jerarquías dentro de las tablas de dimensiones (como Región y Comuna dentro de \texttt{dim\_sede\_geografia}, o Área y Subárea dentro de \texttt{dim\_carrera\_programa}), se maximiza la velocidad de escaneo y compresión columnar del motor VertiPaq en Power BI.

\section{Declaración Formal del Grano}

El grano constituye el cimiento lógico que define el significado exacto de una fila en la tabla de hechos:

\begin{tcolorbox}[colback=lightbg,colframe=secondaryblue,title=\textbf{Declaración Oficial del Grano}]
\emph{«Una fila en \texttt{fact\_matricula\_admision} representa el evento anual de matrícula y postulación para una \textbf{oferta académica específica}, impartida por una \textbf{institución de educación superior}, en una \textbf{sede física o campus de la Región Metropolitana de Santiago}, bajo una \textbf{jornada u horario determinado}, durante un \textbf{año cronológico de admisión (2012 a 2026)}».}
\end{tcolorbox}

\textbf{Propiedades del Grano Atómico Implementado:}
\begin{enumerate}
    \item \textbf{Unicidad Compuesta:} La clave primaria subrogada \texttt{id\_fact} se corresponde biunívocamente con la combinación de las cuatro claves foráneas: $(\texttt{id\_tiempo}, \texttt{id\_sede\_geografia}, \texttt{id\_institucion\_sk}, \texttt{id\_carrera\_programa\_sk})$.
    \item \textbf{Preservación de Variabilidad:} El grano fino evita sesgos de agregación prematura, posibilitando analizar fenómenos de nicho como carreras vespertinas técnicas saturadas frente a planes diurnos profesionales subocupados en la misma sede.
\end{enumerate}

\section{Gobernanza Temporal: Implementación de SCD Tipo 2}

Las instituciones de educación superior en Chile mutan significativamente sus atributos regulatorios a lo largo de un horizonte de 15 años (2012--2026): varían sus años de acreditación CNA, su estatus de autonomía y su adscripción a la política de gratuidad (Ley 20.882 de 2016).

\subsection{Crítica a SCD Tipo 1 (Sobrescritura)}
Si se hubiese aplicado SCD Tipo 1, una institución que en 2012 no contaba con acreditación institucional figuraría retrospectivamente como ``Acreditada con 6 años de Excelencia'' si en 2026 alcanzó dicho reconocimiento. Esto introduce un anacronismo metodológico que invalidaría cualquier correlación entre acreditación y sobredemanda histórica.

\subsection{Mecanismo SCD Tipo 2 Aplicado}
Para mantener fidelidad histórica absoluta:
\begin{itemize}
    \item Se implementó la clave subrogada \texttt{id\_institucion\_sk}, la cual genera un nuevo registro para cada fotografía anual de la institución (totalizando \textbf{1.298 fotografías históricas}).
    \item El atributo \texttt{ano\_vigencia} delimita el marco temporal de validez del estado regulatorio.
    \item En la tabla de hechos, cada oferta histórica se vincula al \texttt{id\_institucion\_sk} del año correspondiente.
    \item Al filtrar por \texttt{ano\_vigencia = 2026}, el modelo refleja con exactitud matemática las \textbf{69 instituciones activas en Santiago}: 35 Universidades, 19 Institutos Profesionales, 11 Centros de Formación Técnica y 4 Instituciones de las FF.AA.
\end{itemize}

\section{Clasificación Kimball de Medidas y Recálculo Dinámico con DIVIDE()}

\subsection{Taxonomía de Medidas (Clase 4 DW Pág. 29)}
\begin{enumerate}
    \item \textbf{Medidas Aditivas Puras (Distributivas):} Operables mediante \texttt{SUM} y \texttt{COUNT} a través de todas las dimensiones. Ejemplos: \texttt{vacantes}, \texttt{matricula\_primer\_ano\_total}, \texttt{matricula\_total}, \texttt{sobrecupo\_neto\_alumnos}, \texttt{secciones\_extra\_por\_sobrecupo}.
    \item \textbf{Medidas Semi-Aditivas:} Sumables a través de geografía, institución o carrera, pero no a lo largo del tiempo (como \texttt{matricula\_total} cuando se analiza como stock de alumnos recurrentes).
    \item \textbf{Medidas No Aditivas (Tasas y Ratios):} Operaciones algebraicas que \textbf{bajo ninguna circunstancia pueden sumarse o promediarse simple}. Ejemplos: \texttt{tasa\_ocupacion\_primer\_ano}, \texttt{tasa\_saturacion\_lab\_35\_pct}, \texttt{puntaje\_promedio\_equiv\_paes}.
\end{enumerate}

\subsection{Demostración Matemática: Falacia de la Agregación de Tasas (Efecto Simpson)}
Consideremos una facultad en Santiago con dos carreras:
\begin{itemize}
    \item \textbf{Carrera A (Plan Común Masivo):} $V_A = 600$ vacantes, $M_A = 660$ matriculados $\implies T_A = 110,0\%$ de ocupación.
    \item \textbf{Carrera B (Especialidad de Nicho):} $V_B = 20$ vacantes, $M_B = 4$ matriculados $\implies T_B = 20,0\%$ de ocupación.
\end{itemize}

Si un analista inexperto promedia aritméticamente las tasas precalculadas:
\[
\text{Tasa Errónea (AVERAGE)} = \frac{110,0\% + 20,0\%}{2} = 65,0\% \quad (\text{Subocupación Aparente del 35\%})
\]
Sin embargo, el cálculo algebraico real basado en los totales aditivos es:
\[
\text{Tasa Real (DIVIDE)} = \frac{\sum M_i}{\sum V_i} \times 100 = \frac{660 + 4}{600 + 20} \times 100 = \frac{664}{620} \times 100 = \mathbf{107,10\%}
\]
El promedio simple subestima la demanda real en más de \textbf{42 puntos porcentuales}, ocultando un sobrecupo que colapsa los laboratorios. Por esta razón técnica, en DAX se formulan las medidas mediante \texttt{DIVIDE(SUM(Num), SUM(Den), BLANK()) * 100}.

\section{Modelado del Caso de Negocio: Capacidad en Laboratorios (Módulo 35)}

El profesor recalcó el impacto de aumentar de \textbf{120 a 140 alumnos} cuando la infraestructura de laboratorios opera en módulos estándar de \textbf{35 puestos}:

\begin{enumerate}
    \item \textbf{Escenario Planificado (120 Vacantes):}
    \[
    \text{Secciones Planificadas} = \lceil 120 / 35 \rceil = 4 \text{ secciones} \implies \text{Capacidad Total} = 4 \times 35 = 140 \text{ puestos.}
    \]
    La saturación es $\frac{120}{140} \times 100 = 85,7\%$, dejando 20 puestos de holgura para repitentes.
    \item \textbf{Escenario de Sobrecupo (140 Matriculados):}
    \[
    \text{Secciones Requeridas} = \lceil 140 / 35 \rceil = 4 \text{ secciones} \implies \text{Saturación} = \frac{140}{140} \times 100 = \mathbf{100,0\%}.
    \]
    La holgura se reduce a cero puestos libres.
    \item \textbf{Quiebre Operacional (Alumno 141):}
    \[
    \text{Secciones Reales} = \lceil 141 / 35 \rceil = 5 \text{ secciones} \implies \text{Secciones Extra} = 5 - 4 = \mathbf{1 \text{ sección extra forzada.}}
    \]
    Un único estudiante adicional obliga a contratar docentes y arrendar horas de laboratorio para 34 puestos vacíos, incrementando los costos fijos directos en un $+25\%$.
\end{enumerate}

\section{Auditoría de Integridad Referencial y Buenas Prácticas DW}

La ejecución del script oficial de auditoría \texttt{audit\_architecture\_and\_integrity\_p3.py} certifica:
\begin{itemize}
    \item \textbf{Integridad Referencial 1:N:} 0 registros huérfanos detectados en las 58.926 filas de la tabla de hechos frente a las dimensiones de tiempo, sede, institución y carrera (\textbf{100\% de integridad}).
    \item \textbf{Cero VARCHAR en Hechos:} Las 43 columnas de \texttt{fact\_matricula\_admision} son 100\% numéricas (\texttt{INTEGER} y \texttt{REAL}), satisfaciendo el mandamiento de Clase 4 DW (Pág. 62).
    \item \textbf{Cero Abreviaturas en Dimensiones:} Se eliminaron los códigos crudos \texttt{Univ.}, \texttt{I.P.}, \texttt{C.F.T.}, \texttt{F.F.A.A.} y \texttt{M/}, reemplazándolos por etiquetas formales.
    \item \textbf{Rendimiento de Consultas:} Consultas analíticas multidimensionales resueltas en menos de \textbf{25 milisegundos}.
\end{itemize}

\end{document}

